\documentclass[10pt, aspectratio=169, compress, protectframetitle, handout]{beamer}
\usepackage[utf8]{inputenc}
\usepackage[english]{babel}
\usepackage{appendixnumberbeamer}
% handout to deactivate \uncover
% usetitleprogressbar might be needed
%\usepackage{beamerprosper}
\usepackage{comment}
% Load BEFORE the theme
\usepackage[normalem]{ulem}
\usepackage[T1]{fontenc}

\usetheme[progressbar=frametitle,block=fill,numbering=fraction]{metropolis}
\setbeamertemplate{blocks}[rounded][shadow=true]
%\setbeamertemplate{note page}[plain]
%\setsansfont[
%     Extension      = .otf,
%     UprightFont    = *-Light,
%     ItalicFont     = *-LightItalic,
%     BoldFont       = *-Regular,
%     BoldItalicFont = *-RegularItalic
% ]{FiraSans}
%\setmonofont[
%     Extension   = .otf,
%     UprightFont = *-Regular,
%     BoldFont    = *-Medium
%]{FiraMono}


\newcommand{\putbg}{\usebackgroundtemplate{\includegraphics[width=\paperwidth,height=\paperheight]{background-vector_169}}}
\newcommand{\putbgdark}{\usebackgroundtemplate{\includegraphics[width=\paperwidth,height=\paperheight]{background-vector-dark_169}}}


\usepackage[export]{adjustbox}
\usepackage[]{enumitem}
\usepackage{datetime}
\usepackage{textpos}
\usepackage{marvosym} % Smile
% Fixes bad positioning of hats
\usefonttheme{professionalfonts}%[onlymath]{serif}
\PassOptionsToPackage{hyphens}{url}\usepackage{hyperref} % to break the links
\hypersetup{
    colorlinks=true,
    linkcolor=blue,  % color of internal links
    urlcolor=blue,   % color of external links
    citecolor=blue,  % color of links to bibliography
}


%%% Bibliografia
\usepackage[autostyle]{csquotes}
\usepackage[backend=biber]{biblatex}
\addbibresource{biblio.bib}


%%% Metadati
\graphicspath{{figures/PNG/}{figures/PDF/}{figures/}}
\newdateformat{monthyear}{\monthname[\THEMONTH] \THEYEAR}
\title{\vspace*{1.5cm}An algorithmic reasoning approach to GNNs}
\subtitle{A project for the \emph{Deep Learning} course}
\author{Angela Carraro, Matteo Scorcia}
\date{}
\institute{\scshape DSSC + IN20 - UniTS
\vfill
\includegraphics[valign=c, height=0.7cm]{logo_dssc_alt}
\hspace*{0.5cm}
\includegraphics[valign=c, height=0.75cm]{Logo_units_blu}
}

\addtobeamertemplate{frametitle}{}{%
\begin{textblock*}{100mm}(0.90\textwidth,-0.94cm)
\includegraphics[valign=c, height=0.4cm]{logo_dssc_alt_white}
\includegraphics[valign=c, height=0.45cm]{Logo_units_white}
\end{textblock*}}

\begin{document}

{\putbg\maketitle}

%\begin{frame}{Contents}
%	\tableofcontents
%\end{frame}

\begin{frame}{Aim of the project}

    \alert{Graph Neural Networks}.
    
\end{frame}

\begin{frame}{Graph Theory}

    Graphs are a general language for describing and modelling complex systems.

    Graphs are an important building block since they can naturally encode an entity-relationship structure, as well as an invariance to permutations (of both nodes and edges) and awareness of input sparsity.
\end{frame}

\begin{frame}{GNN real-world applications}

    GNNs are used for one of three tasks:
    \begin{itemize}[topsep=4pt,itemsep=8pt]
        \item[{$\bullet$}] {\makebox[3.3cm][l]{\alert{\emph{node classification}}:} predict the label of a given node}\\
        \quad $\longrightarrow$ \; E.g., predicting whether a user is a bot in a social network
        
        \item[{$\bullet$}] {\makebox[3.3cm][l]{\alert{\emph{graph classification}}:} make independent predictions specific to each graph}\\
        \quad $\longrightarrow$ \; E.g., property prediction based on molecular graph structures
        
        \item[{$\bullet$}] {\makebox[3.3cm][l]{\alert{\emph{edge prediction}}:} 
        %predict if there exists an edge between two nodes}\\
        infer the edges between nodes in a graph}\\
        \quad $\longrightarrow$ \; E.g., content recommendation in online platforms, predicting drug side-effects
    \end{itemize}

\end{frame}

{\putbgdark
\begin{frame}[standout]
	\begin{center}
		\Large \uncover<+->{Thank you for your attention!}
		
		\Huge\uncover<+->{\Smiley}
	\end{center}
\end{frame}
}

\begin{frame}[allowframebreaks]{References}

	\nocite{*}
	\printbibliography

\end{frame}

\end{document}